%%%%%%%%%%%%%%%%%%%%%%%%%%%%%%%%%%%%%%%%%%%%%%%%%%%%%%%%%%%%
% 5.5 REAL ANALYSIS
% The Composition of Advanced Mathematics
%%%%%%%%%%%%%%%%%%%%%%%%%%%%%%%%%%%%%%%%%%%%%%%%%%%%%%%%%%%%

\section{Real Analysis}
\label{sec:real-analysis}

\prerequisite{
Foundations of mathematical proof, sets and functions, elementary logic,
sequences, limits, continuity, single-variable calculus, and basic
properties of the real numbers.
}

\learningobjectives{
By the end of this section, the reader should be able to:
\begin{itemize}
    \item explain the completeness property of the real numbers;
    \item work with supremum, infimum, maximum, and minimum;
    \item prove convergence or divergence of sequences;
    \item use the Monotone Convergence Theorem and Bolzano--Weierstrass Theorem;
    \item analyze Cauchy sequences;
    \item distinguish pointwise and uniform convergence;
    \item use rigorous epsilon definitions of limits and continuity;
    \item prove basic properties of continuous functions;
    \item apply the Intermediate and Extreme Value Theorems;
    \item distinguish continuity from uniform continuity;
    \item analyze differentiability rigorously;
    \item understand the Mean Value Theorem and its consequences;
    \item define and analyze the Riemann integral;
    \item understand upper and lower sums;
    \item characterize basic Riemann integrability conditions;
    \item analyze sequences and series of functions;
    \item understand power series and uniform convergence;
    \item work with metric-space formulations of fundamental concepts;
    \item connect real analysis to measure theory, functional analysis,
          probability, topology, and differential equations.
\end{itemize}
}

%%%%%%%%%%%%%%%%%%%%%%%%%%%%%%%%%%%%%%%%%%%%%%%%%%%%%%%%%%%%
\subsection{What Is Real Analysis?}
\label{subsec:what-is-real-analysis}
%%%%%%%%%%%%%%%%%%%%%%%%%%%%%%%%%%%%%%%%%%%%%%%%%%%%%%%%%%%%

Calculus develops powerful computational methods involving limits,
derivatives, integrals, and infinite series.

Real analysis studies the logical foundations behind those methods.

Questions such as

\[
\boxed{
\text{What exactly does }x_n\to L\text{ mean?}
}
\]

and

\[
\boxed{
\text{When may a limit pass through an integral?}
}
\]

require precise definitions and proofs.

Real analysis therefore transforms the computational ideas of calculus
into a rigorous mathematical theory.

%%%%%%%%%%%%%%%%%%%%%%%%%%%%%%%%%%%%%%%%%%%%%%%%%%%%%%%%%%%%
\subsection{The Real Number System}
\label{subsec:real-number-system-analysis}
%%%%%%%%%%%%%%%%%%%%%%%%%%%%%%%%%%%%%%%%%%%%%%%%%%%%%%%%%%%%

The real numbers

\[
\R
\]

form an ordered field.

They support the usual arithmetic operations and ordering relations while
satisfying a crucial additional property:

\[
\boxed{
\text{completeness}.
}
\]

Completeness distinguishes \(\R\) from the rational numbers \(\Q\).

%%%%%%%%%%%%%%%%%%%%%%%%%%%%%%%%%%%%%%%%%%%%%%%%%%%%%%%%%%%%
\subsection{Upper and Lower Bounds}
\label{subsec:upper-lower-bounds}
%%%%%%%%%%%%%%%%%%%%%%%%%%%%%%%%%%%%%%%%%%%%%%%%%%%%%%%%%%%%

Let

\[
A\subseteq\R.
\]

A number \(M\in\R\) is an upper bound for \(A\) if

\[
\boxed{
x\leq M
\qquad
\text{for every }x\in A.
}
\]

A number \(m\in\R\) is a lower bound if

\[
\boxed{
m\leq x
\qquad
\text{for every }x\in A.
}
\]

A set having an upper bound is called bounded above.

A set having a lower bound is called bounded below.

%%%%%%%%%%%%%%%%%%%%%%%%%%%%%%%%%%%%%%%%%%%%%%%%%%%%%%%%%%%%
\subsection{Supremum and Infimum}
\label{subsec:supremum-infimum}
%%%%%%%%%%%%%%%%%%%%%%%%%%%%%%%%%%%%%%%%%%%%%%%%%%%%%%%%%%%%

\begin{definitionbox}{Supremum}
If \(A\subseteq\R\) is nonempty and bounded above, its supremum is the
least upper bound of \(A\).

It is denoted

\[
\boxed{
\sup A.
}
\]
\end{definitionbox}

Similarly:

\begin{definitionbox}{Infimum}
If \(A\subseteq\R\) is nonempty and bounded below, its infimum is the
greatest lower bound of \(A\).

It is denoted

\[
\boxed{
\inf A.
}
\]
\end{definitionbox}

The words \emph{supremum} and \emph{infimum} are often abbreviated as
\emph{sup} and \emph{inf}.

%%%%%%%%%%%%%%%%%%%%%%%%%%%%%%%%%%%%%%%%%%%%%%%%%%%%%%%%%%%%
\subsection{Maximum Versus Supremum}
\label{subsec:maximum-versus-supremum}
%%%%%%%%%%%%%%%%%%%%%%%%%%%%%%%%%%%%%%%%%%%%%%%%%%%%%%%%%%%%

A maximum must belong to the set.

A supremum need not.

For example,

\[
A=(0,1).
\]

Then

\[
\boxed{
\sup A=1,
}
\]

but

\[
1\notin A.
\]

Therefore \(A\) has no maximum.

Similarly,

\[
\inf A=0,
\]

but \(A\) has no minimum.

%%%%%%%%%%%%%%%%%%%%%%%%%%%%%%%%%%%%%%%%%%%%%%%%%%%%%%%%%%%%
\subsection{Completeness Axiom}
\label{subsec:completeness-axiom}
%%%%%%%%%%%%%%%%%%%%%%%%%%%%%%%%%%%%%%%%%%%%%%%%%%%%%%%%%%%%

\begin{theorem}[Least Upper Bound Property]
Every nonempty subset of \(\R\) that is bounded above has a supremum in
\(\R\).
\end{theorem}

This property is one standard formulation of the completeness of the real
numbers.

The corresponding greatest lower bound property follows as well.

%%%%%%%%%%%%%%%%%%%%%%%%%%%%%%%%%%%%%%%%%%%%%%%%%%%%%%%%%%%%
\subsection{Why the Rational Numbers Are Not Complete}
\label{subsec:rationals-not-complete}
%%%%%%%%%%%%%%%%%%%%%%%%%%%%%%%%%%%%%%%%%%%%%%%%%%%%%%%%%%%%

Consider

\[
A=
\{q\in\Q:q^2<2\}.
\]

Within \(\R\),

\[
\sup A=\sqrt2.
\]

However,

\[
\sqrt2\notin\Q.
\]

Thus the rational numbers contain a ``gap'' corresponding to this least
upper bound.

The real numbers fill such gaps.

%%%%%%%%%%%%%%%%%%%%%%%%%%%%%%%%%%%%%%%%%%%%%%%%%%%%%%%%%%%%
\subsection{Archimedean Property}
\label{subsec:archimedean-property}
%%%%%%%%%%%%%%%%%%%%%%%%%%%%%%%%%%%%%%%%%%%%%%%%%%%%%%%%%%%%

\begin{theorem}[Archimedean Property]
For every real number \(x\), there exists a natural number \(n\) such that

\[
\boxed{
n>x.
}
\]
\end{theorem}

An important consequence is that for every

\[
\varepsilon>0,
\]

there exists \(n\in\mathbb{N}\) such that

\[
\boxed{
\frac1n<\varepsilon.
}
\]

The lowercase Greek letter

\[
\varepsilon
\]

is called \emph{epsilon}. In analysis it commonly represents an arbitrary
positive tolerance.

%%%%%%%%%%%%%%%%%%%%%%%%%%%%%%%%%%%%%%%%%%%%%%%%%%%%%%%%%%%%
\subsection{Density of the Rational Numbers}
\label{subsec:density-rationals}
%%%%%%%%%%%%%%%%%%%%%%%%%%%%%%%%%%%%%%%%%%%%%%%%%%%%%%%%%%%%

\begin{theorem}
If

\[
a<b,
\]

then there exists

\[
q\in\Q
\]

such that

\[
\boxed{
a<q<b.
}
\]
\end{theorem}

Thus the rational numbers are dense in the real numbers.

The irrational numbers are also dense in \(\R\).

%%%%%%%%%%%%%%%%%%%%%%%%%%%%%%%%%%%%%%%%%%%%%%%%%%%%%%%%%%%%
\subsection{Sequences}
\label{subsec:sequences-real-analysis}
%%%%%%%%%%%%%%%%%%%%%%%%%%%%%%%%%%%%%%%%%%%%%%%%%%%%%%%%%%%%

A sequence is a function

\[
\boxed{
x:\mathbb{N}\to\R.
}
\]

We usually write

\[
x_n
\]

instead of

\[
x(n).
\]

A sequence is therefore an ordered list

\[
x_1,x_2,x_3,\ldots.
\]

%%%%%%%%%%%%%%%%%%%%%%%%%%%%%%%%%%%%%%%%%%%%%%%%%%%%%%%%%%%%
\subsection{Convergence of Sequences}
\label{subsec:sequence-convergence-rigorous}
%%%%%%%%%%%%%%%%%%%%%%%%%%%%%%%%%%%%%%%%%%%%%%%%%%%%%%%%%%%%

\begin{definitionbox}{Convergent Sequence}
A sequence \((x_n)\) converges to \(L\in\R\) if for every
\(\varepsilon>0\), there exists \(N\in\mathbb{N}\) such that

\[
\boxed{
n\geq N
\Longrightarrow
|x_n-L|<\varepsilon.
}
\]

We write

\[
\boxed{
x_n\to L
}
\]

or

\[
\boxed{
\lim_{n\to\infty}x_n=L.
}
\]
\end{definitionbox}

The definition says that the terms eventually remain inside every
prescribed neighborhood of \(L\).

%%%%%%%%%%%%%%%%%%%%%%%%%%%%%%%%%%%%%%%%%%%%%%%%%%%%%%%%%%%%
\subsection{Worked Example: Epsilon Proof for a Sequence}
\label{subsec:worked-epsilon-sequence}
%%%%%%%%%%%%%%%%%%%%%%%%%%%%%%%%%%%%%%%%%%%%%%%%%%%%%%%%%%%%

\begin{workedexamplebox}{Proving \(1/n\to0\)}

We prove

\[
\lim_{n\to\infty}\frac1n=0.
\]

Let

\[
\varepsilon>0.
\]

We want

\[
\left|
\frac1n-0
\right|
<
\varepsilon.
\]

This is equivalent to

\[
\frac1n<\varepsilon,
\]

or

\[
n>\frac1\varepsilon.
\]

By the Archimedean Property, choose

\[
N>\frac1\varepsilon.
\]

Then whenever

\[
n\geq N,
\]

we have

\[
\left|
\frac1n
\right|
\leq
\frac1N
<
\varepsilon.
\]

\begin{answerbox}
\[
\boxed{
\frac1n\to0.
}
\]
\end{answerbox}
\end{workedexamplebox}

%%%%%%%%%%%%%%%%%%%%%%%%%%%%%%%%%%%%%%%%%%%%%%%%%%%%%%%%%%%%
\subsection{Uniqueness of Limits}
\label{subsec:uniqueness-sequence-limits}
%%%%%%%%%%%%%%%%%%%%%%%%%%%%%%%%%%%%%%%%%%%%%%%%%%%%%%%%%%%%

\begin{theorem}[Uniqueness of Limits]
A convergent sequence has exactly one limit.
\end{theorem}

\begin{proof}
Suppose

\[
x_n\to L
\]

and

\[
x_n\to M.
\]

Assume

\[
L\neq M.
\]

Choose

\[
\varepsilon
=
\frac{|L-M|}{3}.
\]

For sufficiently large \(n\),

\[
|x_n-L|<\varepsilon
\]

and

\[
|x_n-M|<\varepsilon.
\]

By the triangle inequality,

\[
|L-M|
\leq
|L-x_n|+|x_n-M|
<
2\varepsilon
=
\frac23|L-M|,
\]

which is impossible.

Therefore

\[
L=M.
\]
\end{proof}

%%%%%%%%%%%%%%%%%%%%%%%%%%%%%%%%%%%%%%%%%%%%%%%%%%%%%%%%%%%%
\subsection{Bounded Sequences}
\label{subsec:bounded-sequences}
%%%%%%%%%%%%%%%%%%%%%%%%%%%%%%%%%%%%%%%%%%%%%%%%%%%%%%%%%%%%

A sequence \((x_n)\) is bounded if there exists

\[
M>0
\]

such that

\[
\boxed{
|x_n|\leq M
}
\]

for every \(n\).

\begin{theorem}
Every convergent real sequence is bounded.
\end{theorem}

The converse is false.

For example,

\[
x_n=(-1)^n
\]

is bounded but does not converge.

%%%%%%%%%%%%%%%%%%%%%%%%%%%%%%%%%%%%%%%%%%%%%%%%%%%%%%%%%%%%
\subsection{Algebra of Sequence Limits}
\label{subsec:algebra-sequence-limits}
%%%%%%%%%%%%%%%%%%%%%%%%%%%%%%%%%%%%%%%%%%%%%%%%%%%%%%%%%%%%

If

\[
x_n\to x
\]

and

\[
y_n\to y,
\]

then

\[
\boxed{
x_n+y_n\to x+y,
}
\]

\[
\boxed{
x_ny_n\to xy,
}
\]

and for any constant \(c\),

\[
\boxed{
cx_n\to cx.
}
\]

If

\[
y\neq0
\]

and \(y_n\neq0\) eventually, then

\[
\boxed{
\frac{x_n}{y_n}
\to
\frac{x}{y}.
}
\]

%%%%%%%%%%%%%%%%%%%%%%%%%%%%%%%%%%%%%%%%%%%%%%%%%%%%%%%%%%%%
\subsection{Monotone Sequences}
\label{subsec:monotone-sequences}
%%%%%%%%%%%%%%%%%%%%%%%%%%%%%%%%%%%%%%%%%%%%%%%%%%%%%%%%%%%%

A sequence is increasing if

\[
\boxed{
x_{n+1}\geq x_n
}
\]

for every \(n\).

It is decreasing if

\[
\boxed{
x_{n+1}\leq x_n.
}
\]

A sequence that is either increasing or decreasing is called monotone.

%%%%%%%%%%%%%%%%%%%%%%%%%%%%%%%%%%%%%%%%%%%%%%%%%%%%%%%%%%%%
\subsection{Monotone Convergence Theorem}
\label{subsec:monotone-convergence-sequences}
%%%%%%%%%%%%%%%%%%%%%%%%%%%%%%%%%%%%%%%%%%%%%%%%%%%%%%%%%%%%

\begin{theorem}[Monotone Convergence Theorem for Sequences]
Every increasing sequence that is bounded above converges.

Every decreasing sequence that is bounded below converges.
\end{theorem}

For an increasing bounded sequence,

\[
\boxed{
\lim_{n\to\infty}x_n
=
\sup\{x_n:n\in\mathbb{N}\}.
}
\]

This theorem is a direct manifestation of completeness.

%%%%%%%%%%%%%%%%%%%%%%%%%%%%%%%%%%%%%%%%%%%%%%%%%%%%%%%%%%%%
\subsection{Subsequences}
\label{subsec:subsequences}
%%%%%%%%%%%%%%%%%%%%%%%%%%%%%%%%%%%%%%%%%%%%%%%%%%%%%%%%%%%%

A subsequence is obtained by selecting terms whose indices form a strictly
increasing sequence

\[
n_1<n_2<n_3<\cdots.
\]

The resulting subsequence is written

\[
\boxed{
(x_{n_k}).
}
\]

If

\[
x_n\to L,
\]

then every subsequence also converges to \(L\).

%%%%%%%%%%%%%%%%%%%%%%%%%%%%%%%%%%%%%%%%%%%%%%%%%%%%%%%%%%%%
\subsection{Bolzano--Weierstrass Theorem}
\label{subsec:bolzano-weierstrass}
%%%%%%%%%%%%%%%%%%%%%%%%%%%%%%%%%%%%%%%%%%%%%%%%%%%%%%%%%%%%

\begin{theorem}[Bolzano--Weierstrass Theorem]
Every bounded sequence in \(\R\) has a convergent subsequence.
\end{theorem}

More generally, every bounded sequence in

\[
\R^n
\]

has a convergent subsequence.

This theorem is closely related to compactness.

%%%%%%%%%%%%%%%%%%%%%%%%%%%%%%%%%%%%%%%%%%%%%%%%%%%%%%%%%%%%
\subsection{Limit Superior and Limit Inferior}
\label{subsec:limsup-liminf}
%%%%%%%%%%%%%%%%%%%%%%%%%%%%%%%%%%%%%%%%%%%%%%%%%%%%%%%%%%%%

For a bounded sequence \((x_n)\), define

\[
s_n
=
\sup\{x_k:k\geq n\}
\]

and

\[
i_n
=
\inf\{x_k:k\geq n\}.
\]

Then the limit superior is

\[
\boxed{
\limsup_{n\to\infty}x_n
=
\lim_{n\to\infty}s_n,
}
\]

and the limit inferior is

\[
\boxed{
\liminf_{n\to\infty}x_n
=
\lim_{n\to\infty}i_n.
}
\]

These describe the eventual upper and lower limiting behavior of a
sequence.

%%%%%%%%%%%%%%%%%%%%%%%%%%%%%%%%%%%%%%%%%%%%%%%%%%%%%%%%%%%%
\subsection{Cauchy Sequences}
\label{subsec:cauchy-sequences}
%%%%%%%%%%%%%%%%%%%%%%%%%%%%%%%%%%%%%%%%%%%%%%%%%%%%%%%%%%%%

\begin{definitionbox}{Cauchy Sequence}
A sequence \((x_n)\) is Cauchy if for every \(\varepsilon>0\), there exists
\(N\in\mathbb{N}\) such that

\[
\boxed{
m,n\geq N
\Longrightarrow
|x_m-x_n|<\varepsilon.
}
\]
\end{definitionbox}

The definition does not mention a candidate limit.

It states that sufficiently late terms become arbitrarily close to one
another.

%%%%%%%%%%%%%%%%%%%%%%%%%%%%%%%%%%%%%%%%%%%%%%%%%%%%%%%%%%%%
\subsection{Cauchy Criterion for the Real Numbers}
\label{subsec:cauchy-criterion-real}
%%%%%%%%%%%%%%%%%%%%%%%%%%%%%%%%%%%%%%%%%%%%%%%%%%%%%%%%%%%%

\begin{theorem}[Cauchy Criterion]
A sequence of real numbers converges if and only if it is Cauchy.
\end{theorem}

The implication

\[
\text{convergent}
\Longrightarrow
\text{Cauchy}
\]

holds in every metric space.

The reverse implication depends on completeness.

%%%%%%%%%%%%%%%%%%%%%%%%%%%%%%%%%%%%%%%%%%%%%%%%%%%%%%%%%%%%
\subsection{Series}
\label{subsec:series-real-analysis}
%%%%%%%%%%%%%%%%%%%%%%%%%%%%%%%%%%%%%%%%%%%%%%%%%%%%%%%%%%%%

An infinite series

\[
\sum_{n=1}^{\infty}a_n
\]

is defined through its sequence of partial sums

\[
\boxed{
s_N
=
\sum_{n=1}^{N}a_n.
}
\]

The series converges if the sequence

\[
(s_N)
\]

converges.

Thus series theory is fundamentally sequence theory.

%%%%%%%%%%%%%%%%%%%%%%%%%%%%%%%%%%%%%%%%%%%%%%%%%%%%%%%%%%%%
\subsection{Necessary Condition for Series Convergence}
\label{subsec:necessary-series-condition}
%%%%%%%%%%%%%%%%%%%%%%%%%%%%%%%%%%%%%%%%%%%%%%%%%%%%%%%%%%%%

\begin{theorem}
If

\[
\sum_{n=1}^{\infty}a_n
\]

converges, then

\[
\boxed{
a_n\to0.
}
\]
\end{theorem}

The converse is false.

The harmonic series

\[
\sum_{n=1}^{\infty}\frac1n
\]

diverges even though

\[
\frac1n\to0.
\]

%%%%%%%%%%%%%%%%%%%%%%%%%%%%%%%%%%%%%%%%%%%%%%%%%%%%%%%%%%%%
\subsection{Absolute Convergence}
\label{subsec:absolute-convergence-analysis}
%%%%%%%%%%%%%%%%%%%%%%%%%%%%%%%%%%%%%%%%%%%%%%%%%%%%%%%%%%%%

A series

\[
\sum a_n
\]

converges absolutely if

\[
\boxed{
\sum |a_n|
}
\]

converges.

\begin{theorem}
Every absolutely convergent real series converges.
\end{theorem}

A convergent series that is not absolutely convergent is called
conditionally convergent.

%%%%%%%%%%%%%%%%%%%%%%%%%%%%%%%%%%%%%%%%%%%%%%%%%%%%%%%%%%%%
\subsection{Rearrangements of Series}
\label{subsec:rearrangements-series}
%%%%%%%%%%%%%%%%%%%%%%%%%%%%%%%%%%%%%%%%%%%%%%%%%%%%%%%%%%%%

Absolutely convergent series may be rearranged without changing their sum.

Conditionally convergent series behave very differently.

The Riemann rearrangement theorem shows that rearrangements of a
conditionally convergent real series can dramatically alter its behavior.

This illustrates why absolute convergence is a much stronger property than
ordinary convergence.

%%%%%%%%%%%%%%%%%%%%%%%%%%%%%%%%%%%%%%%%%%%%%%%%%%%%%%%%%%%%
\subsection{Limits of Functions}
\label{subsec:function-limits-rigorous}
%%%%%%%%%%%%%%%%%%%%%%%%%%%%%%%%%%%%%%%%%%%%%%%%%%%%%%%%%%%%

\begin{definitionbox}{Limit of a Function}
Let \(f\) be defined near \(a\), possibly excluding \(a\). We say

\[
\lim_{x\to a}f(x)=L
\]

if for every

\[
\varepsilon>0,
\]

there exists

\[
\delta>0
\]

such that

\[
\boxed{
0<|x-a|<\delta
\Longrightarrow
|f(x)-L|<\varepsilon.
}
\]

The lowercase Greek letter \(\delta\), read \emph{delta}, represents a
positive restriction on the input.
\end{definitionbox}

The definition may be summarized as

\[
\boxed{
\forall\varepsilon>0
\quad
\exists\delta>0
\quad
0<|x-a|<\delta
\Rightarrow
|f(x)-L|<\varepsilon.
}
\]

%%%%%%%%%%%%%%%%%%%%%%%%%%%%%%%%%%%%%%%%%%%%%%%%%%%%%%%%%%%%
\subsection{Worked Example: Epsilon--Delta Proof}
\label{subsec:worked-epsilon-delta}
%%%%%%%%%%%%%%%%%%%%%%%%%%%%%%%%%%%%%%%%%%%%%%%%%%%%%%%%%%%%

\begin{workedexamplebox}{A Linear Function}

Prove

\[
\lim_{x\to2}(3x+1)=7.
\]

Let

\[
\varepsilon>0.
\]

We need

\[
|(3x+1)-7|<\varepsilon.
\]

Simplifying,

\[
|(3x+1)-7|
=
|3x-6|
=
3|x-2|.
\]

Choose

\[
\delta=\frac{\varepsilon}{3}.
\]

Then

\[
0<|x-2|<\delta
\]

implies

\[
|(3x+1)-7|
=
3|x-2|
<
3\delta
=
\varepsilon.
\]

\begin{answerbox}
\[
\boxed{
\lim_{x\to2}(3x+1)=7.
}
\]
\end{answerbox}
\end{workedexamplebox}

%%%%%%%%%%%%%%%%%%%%%%%%%%%%%%%%%%%%%%%%%%%%%%%%%%%%%%%%%%%%
\subsection{Sequential Characterization of Limits}
\label{subsec:sequential-characterization-limits}
%%%%%%%%%%%%%%%%%%%%%%%%%%%%%%%%%%%%%%%%%%%%%%%%%%%%%%%%%%%%

\begin{theorem}[Sequential Criterion]
The statement

\[
\lim_{x\to a}f(x)=L
\]

holds if and only if for every sequence \((x_n)\) satisfying

\[
x_n\neq a
\]

and

\[
x_n\to a,
\]

we have

\[
\boxed{
f(x_n)\to L.
}
\]
\end{theorem}

This theorem connects function limits directly to sequence convergence.

%%%%%%%%%%%%%%%%%%%%%%%%%%%%%%%%%%%%%%%%%%%%%%%%%%%%%%%%%%%%
\subsection{Continuity}
\label{subsec:continuity-real-analysis}
%%%%%%%%%%%%%%%%%%%%%%%%%%%%%%%%%%%%%%%%%%%%%%%%%%%%%%%%%%%%

\begin{definitionbox}{Continuity}
A function \(f\) is continuous at \(a\) if for every \(\varepsilon>0\),
there exists \(\delta>0\) such that

\[
\boxed{
|x-a|<\delta
\Longrightarrow
|f(x)-f(a)|<\varepsilon.
}
\]
\end{definitionbox}

Equivalently,

\[
\boxed{
\lim_{x\to a}f(x)=f(a).
}
\]

A function is continuous on a set if it is continuous at every point of
that set, with the appropriate relative-domain interpretation at boundary
points.

%%%%%%%%%%%%%%%%%%%%%%%%%%%%%%%%%%%%%%%%%%%%%%%%%%%%%%%%%%%%
\subsection{Sequential Characterization of Continuity}
\label{subsec:sequential-continuity}
%%%%%%%%%%%%%%%%%%%%%%%%%%%%%%%%%%%%%%%%%%%%%%%%%%%%%%%%%%%%

A function \(f\) is continuous at \(a\) if and only if

\[
\boxed{
x_n\to a
\Longrightarrow
f(x_n)\to f(a).
}
\]

This formulation is often especially useful in proofs.

%%%%%%%%%%%%%%%%%%%%%%%%%%%%%%%%%%%%%%%%%%%%%%%%%%%%%%%%%%%%
\subsection{Algebra of Continuous Functions}
\label{subsec:algebra-continuous-functions}
%%%%%%%%%%%%%%%%%%%%%%%%%%%%%%%%%%%%%%%%%%%%%%%%%%%%%%%%%%%%

If \(f\) and \(g\) are continuous at \(a\), then so are

\[
\boxed{
f+g,
\qquad
fg,
\qquad
cf.
}
\]

If

\[
g(a)\neq0,
\]

then

\[
\boxed{
\frac{f}{g}
}
\]

is continuous at \(a\).

Compositions of continuous functions are continuous wherever the
composition is defined.

%%%%%%%%%%%%%%%%%%%%%%%%%%%%%%%%%%%%%%%%%%%%%%%%%%%%%%%%%%%%
\subsection{Intermediate Value Theorem}
\label{subsec:intermediate-value-theorem-analysis}
%%%%%%%%%%%%%%%%%%%%%%%%%%%%%%%%%%%%%%%%%%%%%%%%%%%%%%%%%%%%

\begin{theorem}[Intermediate Value Theorem]
Suppose

\[
f:[a,b]\to\R
\]

is continuous.

If \(N\) lies between \(f(a)\) and \(f(b)\), then there exists

\[
c\in[a,b]
\]

such that

\[
\boxed{
f(c)=N.
}
\]
\end{theorem}

A continuous function on an interval cannot jump over intermediate values.

%%%%%%%%%%%%%%%%%%%%%%%%%%%%%%%%%%%%%%%%%%%%%%%%%%%%%%%%%%%%
\subsection{Existence of Roots}
\label{subsec:existence-roots-ivt}
%%%%%%%%%%%%%%%%%%%%%%%%%%%%%%%%%%%%%%%%%%%%%%%%%%%%%%%%%%%%

A particularly useful consequence occurs when

\[
f(a)f(b)<0.
\]

Then \(f(a)\) and \(f(b)\) have opposite signs.

By the Intermediate Value Theorem, there exists

\[
c\in(a,b)
\]

such that

\[
\boxed{
f(c)=0.
}
\]

This gives an important existence theorem even when the root cannot be
written explicitly.

%%%%%%%%%%%%%%%%%%%%%%%%%%%%%%%%%%%%%%%%%%%%%%%%%%%%%%%%%%%%
\subsection{Extreme Value Theorem}
\label{subsec:extreme-value-theorem-analysis}
%%%%%%%%%%%%%%%%%%%%%%%%%%%%%%%%%%%%%%%%%%%%%%%%%%%%%%%%%%%%

\begin{theorem}[Extreme Value Theorem]
If

\[
f:[a,b]\to\R
\]

is continuous, then \(f\) attains both an absolute maximum and an absolute
minimum on \([a,b]\).
\end{theorem}

Thus there exist points

\[
x_{\min},x_{\max}\in[a,b]
\]

such that

\[
\boxed{
f(x_{\min})
\leq
f(x)
\leq
f(x_{\max})
}
\]

for every \(x\in[a,b]\).

%%%%%%%%%%%%%%%%%%%%%%%%%%%%%%%%%%%%%%%%%%%%%%%%%%%%%%%%%%%%
\subsection{Compactness and Continuity}
\label{subsec:compactness-continuity-analysis}
%%%%%%%%%%%%%%%%%%%%%%%%%%%%%%%%%%%%%%%%%%%%%%%%%%%%%%%%%%%%

The Extreme Value Theorem is part of a broader principle:

\[
\boxed{
\text{continuous images of compact sets are compact}.
}
\]

Since compact subsets of \(\R\) are closed and bounded, the image

\[
f([a,b])
\]

is also compact and therefore contains its supremum and infimum.

This reveals the topological structure behind the theorem.

%%%%%%%%%%%%%%%%%%%%%%%%%%%%%%%%%%%%%%%%%%%%%%%%%%%%%%%%%%%%
\subsection{Uniform Continuity}
\label{subsec:uniform-continuity}
%%%%%%%%%%%%%%%%%%%%%%%%%%%%%%%%%%%%%%%%%%%%%%%%%%%%%%%%%%%%

\begin{definitionbox}{Uniform Continuity}
A function

\[
f:A\to\R
\]

is uniformly continuous on \(A\) if for every \(\varepsilon>0\), there
exists \(\delta>0\) such that for every \(x,y\in A\),

\[
\boxed{
|x-y|<\delta
\Longrightarrow
|f(x)-f(y)|<\varepsilon.
}
\]
\end{definitionbox}

The important distinction is that the same \(\delta\) works throughout the
entire domain.

%%%%%%%%%%%%%%%%%%%%%%%%%%%%%%%%%%%%%%%%%%%%%%%%%%%%%%%%%%%%
\subsection{Continuity Versus Uniform Continuity}
\label{subsec:continuity-vs-uniform}
%%%%%%%%%%%%%%%%%%%%%%%%%%%%%%%%%%%%%%%%%%%%%%%%%%%%%%%%%%%%

Ordinary continuity allows

\[
\delta
\]

to depend on both

\[
\varepsilon
\]

and the point.

Uniform continuity requires

\[
\boxed{
\delta=\delta(\varepsilon)
}
\]

independent of the point.

For example,

\[
f(x)=x^2
\]

is continuous on \(\R\), but it is not uniformly continuous on \(\R\).

%%%%%%%%%%%%%%%%%%%%%%%%%%%%%%%%%%%%%%%%%%%%%%%%%%%%%%%%%%%%
\subsection{Heine--Cantor Theorem}
\label{subsec:heine-cantor}
%%%%%%%%%%%%%%%%%%%%%%%%%%%%%%%%%%%%%%%%%%%%%%%%%%%%%%%%%%%%

\begin{theorem}[Heine--Cantor Theorem]
Every continuous function on a compact metric space is uniformly
continuous.
\end{theorem}

In particular, every continuous function

\[
f:[a,b]\to\R
\]

is uniformly continuous.

%%%%%%%%%%%%%%%%%%%%%%%%%%%%%%%%%%%%%%%%%%%%%%%%%%%%%%%%%%%%
\subsection{Differentiability}
\label{subsec:differentiability-real-analysis}
%%%%%%%%%%%%%%%%%%%%%%%%%%%%%%%%%%%%%%%%%%%%%%%%%%%%%%%%%%%%

A function \(f\) is differentiable at \(a\) if the limit

\[
\boxed{
f'(a)
=
\lim_{h\to0}
\frac{f(a+h)-f(a)}{h}
}
\]

exists as a finite real number.

Equivalently,

\[
\boxed{
f(a+h)
=
f(a)
+
f'(a)h
+
o(h).
}
\]

The notation

\[
o(h)
\]

means an error term satisfying

\[
\frac{o(h)}{h}\to0
\]

as \(h\to0\).

%%%%%%%%%%%%%%%%%%%%%%%%%%%%%%%%%%%%%%%%%%%%%%%%%%%%%%%%%%%%
\subsection{Differentiability Implies Continuity}
\label{subsec:differentiability-implies-continuity}
%%%%%%%%%%%%%%%%%%%%%%%%%%%%%%%%%%%%%%%%%%%%%%%%%%%%%%%%%%%%

\begin{theorem}
If \(f\) is differentiable at \(a\), then \(f\) is continuous at \(a\).
\end{theorem}

\begin{proof}
For \(x\neq a\),

\[
f(x)-f(a)
=
\frac{f(x)-f(a)}{x-a}(x-a).
\]

As

\[
x\to a,
\]

the first factor approaches

\[
f'(a)
\]

and the second approaches \(0\).

Therefore

\[
f(x)-f(a)\to0,
\]

so

\[
f(x)\to f(a).
\]

Hence \(f\) is continuous at \(a\).
\end{proof}

The converse is false.

For example,

\[
f(x)=|x|
\]

is continuous but not differentiable at \(0\).

%%%%%%%%%%%%%%%%%%%%%%%%%%%%%%%%%%%%%%%%%%%%%%%%%%%%%%%%%%%%
\subsection{Rolle's Theorem}
\label{subsec:rolles-theorem}
%%%%%%%%%%%%%%%%%%%%%%%%%%%%%%%%%%%%%%%%%%%%%%%%%%%%%%%%%%%%

\begin{theorem}[Rolle's Theorem]
Suppose \(f\) is continuous on \([a,b]\), differentiable on \((a,b)\), and

\[
f(a)=f(b).
\]

Then there exists

\[
c\in(a,b)
\]

such that

\[
\boxed{
f'(c)=0.
}
\]
\end{theorem}

%%%%%%%%%%%%%%%%%%%%%%%%%%%%%%%%%%%%%%%%%%%%%%%%%%%%%%%%%%%%
\subsection{Mean Value Theorem}
\label{subsec:mean-value-theorem-analysis}
%%%%%%%%%%%%%%%%%%%%%%%%%%%%%%%%%%%%%%%%%%%%%%%%%%%%%%%%%%%%

\begin{theorem}[Mean Value Theorem]
Suppose \(f\) is continuous on \([a,b]\) and differentiable on \((a,b)\).

Then there exists

\[
c\in(a,b)
\]

such that

\[
\boxed{
f'(c)
=
\frac{f(b)-f(a)}{b-a}.
}
\]
\end{theorem}

The theorem states that some instantaneous rate of change equals the
average rate of change over the interval.

%%%%%%%%%%%%%%%%%%%%%%%%%%%%%%%%%%%%%%%%%%%%%%%%%%%%%%%%%%%%
\subsection{Consequences of the Mean Value Theorem}
\label{subsec:mvt-consequences}
%%%%%%%%%%%%%%%%%%%%%%%%%%%%%%%%%%%%%%%%%%%%%%%%%%%%%%%%%%%%

If

\[
f'(x)=0
\]

throughout an interval, then \(f\) is constant there.

If

\[
f'(x)>0,
\]

then \(f\) is strictly increasing.

If

\[
f'(x)<0,
\]

then \(f\) is strictly decreasing.

If

\[
|f'(x)|\leq M,
\]

then the Mean Value Theorem gives

\[
\boxed{
|f(x)-f(y)|
\leq
M|x-y|.
}
\]

This is a Lipschitz estimate.

%%%%%%%%%%%%%%%%%%%%%%%%%%%%%%%%%%%%%%%%%%%%%%%%%%%%%%%%%%%%
\subsection{Lipschitz Continuity}
\label{subsec:lipschitz-continuity}
%%%%%%%%%%%%%%%%%%%%%%%%%%%%%%%%%%%%%%%%%%%%%%%%%%%%%%%%%%%%

\begin{definitionbox}{Lipschitz Continuity}
A function \(f:A\to\R\) is Lipschitz continuous if there exists \(L\geq0\)
such that

\[
\boxed{
|f(x)-f(y)|
\leq
L|x-y|
}
\]

for every \(x,y\in A\).
\end{definitionbox}

Every Lipschitz function is uniformly continuous.

Thus

\[
\boxed{
\text{Lipschitz}
\Longrightarrow
\text{uniformly continuous}
\Longrightarrow
\text{continuous}.
}
\]

The converses do not hold in general.

%%%%%%%%%%%%%%%%%%%%%%%%%%%%%%%%%%%%%%%%%%%%%%%%%%%%%%%%%%%%
\subsection{Taylor's Theorem}
\label{subsec:taylor-theorem-analysis}
%%%%%%%%%%%%%%%%%%%%%%%%%%%%%%%%%%%%%%%%%%%%%%%%%%%%%%%%%%%%

For a sufficiently differentiable function, Taylor's theorem approximates
the function by a polynomial and controls the error.

One common form is

\[
\boxed{
f(x)
=
\sum_{k=0}^{n}
\frac{f^{(k)}(a)}{k!}(x-a)^k
+
R_n(x).
}
\]

Under suitable hypotheses, the Lagrange remainder is

\[
\boxed{
R_n(x)
=
\frac{f^{(n+1)}(\xi)}
{(n+1)!}
(x-a)^{n+1}
}
\]

for some \(\xi\) between \(a\) and \(x\).

Taylor's theorem is therefore both an approximation theorem and an error
estimate.

%%%%%%%%%%%%%%%%%%%%%%%%%%%%%%%%%%%%%%%%%%%%%%%%%%%%%%%%%%%%
\subsection{Partitions of an Interval}
\label{subsec:partitions-riemann}
%%%%%%%%%%%%%%%%%%%%%%%%%%%%%%%%%%%%%%%%%%%%%%%%%%%%%%%%%%%%

Let

\[
[a,b]
\]

be a closed interval.

A partition \(P\) is a finite ordered set

\[
\boxed{
P=
\{x_0,x_1,\ldots,x_n\}
}
\]

satisfying

\[
a=x_0<x_1<\cdots<x_n=b.
\]

The subinterval lengths are

\[
\boxed{
\Delta x_i
=
x_i-x_{i-1}.
}
\]

%%%%%%%%%%%%%%%%%%%%%%%%%%%%%%%%%%%%%%%%%%%%%%%%%%%%%%%%%%%%
\subsection{Upper and Lower Sums}
\label{subsec:upper-lower-riemann-sums}
%%%%%%%%%%%%%%%%%%%%%%%%%%%%%%%%%%%%%%%%%%%%%%%%%%%%%%%%%%%%

Suppose \(f\) is bounded on \([a,b]\).

For each subinterval

\[
[x_{i-1},x_i],
\]

define

\[
M_i
=
\sup
\{f(x):x\in[x_{i-1},x_i]\}
\]

and

\[
m_i
=
\inf
\{f(x):x\in[x_{i-1},x_i]\}.
\]

The upper sum is

\[
\boxed{
U(f,P)
=
\sum_{i=1}^{n}
M_i\Delta x_i,
}
\]

and the lower sum is

\[
\boxed{
L(f,P)
=
\sum_{i=1}^{n}
m_i\Delta x_i.
}
\]

%%%%%%%%%%%%%%%%%%%%%%%%%%%%%%%%%%%%%%%%%%%%%%%%%%%%%%%%%%%%
\subsection{Riemann Integrability}
\label{subsec:riemann-integrability}
%%%%%%%%%%%%%%%%%%%%%%%%%%%%%%%%%%%%%%%%%%%%%%%%%%%%%%%%%%%%

\begin{definitionbox}{Riemann Integrable Function}
A bounded function \(f:[a,b]\to\R\) is Riemann integrable if its upper and
lower integrals agree.

Equivalently, \(f\) is Riemann integrable if for every
\(\varepsilon>0\), there exists a partition \(P\) such that

\[
\boxed{
U(f,P)-L(f,P)<\varepsilon.
}
\]
\end{definitionbox}

The common value is written

\[
\boxed{
\int_a^b f(x)\,dx.
}
\]

%%%%%%%%%%%%%%%%%%%%%%%%%%%%%%%%%%%%%%%%%%%%%%%%%%%%%%%%%%%%
\subsection{Continuous Functions Are Riemann Integrable}
\label{subsec:continuous-riemann-integrable}
%%%%%%%%%%%%%%%%%%%%%%%%%%%%%%%%%%%%%%%%%%%%%%%%%%%%%%%%%%%%

\begin{theorem}
Every continuous function on a closed bounded interval

\[
[a,b]
\]

is Riemann integrable.
\end{theorem}

A key reason is that continuity on a compact interval implies uniform
continuity.

Thus sufficiently small subintervals force the oscillation of \(f\) to be
small.

%%%%%%%%%%%%%%%%%%%%%%%%%%%%%%%%%%%%%%%%%%%%%%%%%%%%%%%%%%%%
\subsection{Monotone Functions Are Integrable}
\label{subsec:monotone-riemann-integrable}
%%%%%%%%%%%%%%%%%%%%%%%%%%%%%%%%%%%%%%%%%%%%%%%%%%%%%%%%%%%%

\begin{theorem}
Every bounded monotone function on a closed bounded interval is Riemann
integrable.
\end{theorem}

Therefore continuity is sufficient for Riemann integrability but is not
necessary.

%%%%%%%%%%%%%%%%%%%%%%%%%%%%%%%%%%%%%%%%%%%%%%%%%%%%%%%%%%%%
\subsection{Discontinuous Functions and Integrability}
\label{subsec:discontinuous-riemann-functions}
%%%%%%%%%%%%%%%%%%%%%%%%%%%%%%%%%%%%%%%%%%%%%%%%%%%%%%%%%%%%

Some discontinuous functions are Riemann integrable.

For example, a bounded function with finitely many discontinuities on a
closed interval is Riemann integrable.

However, not every bounded function is Riemann integrable.

The Dirichlet function

\[
f(x)
=
\begin{cases}
1, & x\in\Q,\\
0, & x\notin\Q,
\end{cases}
\]

is discontinuous everywhere and is not Riemann integrable on any
nondegenerate closed interval.

%%%%%%%%%%%%%%%%%%%%%%%%%%%%%%%%%%%%%%%%%%%%%%%%%%%%%%%%%%%%
\subsection{Fundamental Theorem of Calculus Revisited}
\label{subsec:ftc-real-analysis}
%%%%%%%%%%%%%%%%%%%%%%%%%%%%%%%%%%%%%%%%%%%%%%%%%%%%%%%%%%%%

The Fundamental Theorem of Calculus connects differentiation and
integration.

One form states that if \(f\) is continuous on \([a,b]\) and

\[
F(x)
=
\int_a^x f(t)\,dt,
\]

then

\[
\boxed{
F'(x)=f(x)
}
\]

for

\[
x\in(a,b).
\]

Another form states that if \(F'=f\), then

\[
\boxed{
\int_a^b f(x)\,dx
=
F(b)-F(a).
}
\]

%%%%%%%%%%%%%%%%%%%%%%%%%%%%%%%%%%%%%%%%%%%%%%%%%%%%%%%%%%%%
\subsection{Improper Integrals}
\label{subsec:improper-integrals-analysis}
%%%%%%%%%%%%%%%%%%%%%%%%%%%%%%%%%%%%%%%%%%%%%%%%%%%%%%%%%%%%

An improper integral is defined through limits.

For example,

\[
\boxed{
\int_a^\infty f(x)\,dx
=
\lim_{b\to\infty}
\int_a^b f(x)\,dx
}
\]

when the limit exists.

Similarly, if \(f\) becomes unbounded near \(a\),

\[
\boxed{
\int_a^b f(x)\,dx
=
\lim_{c\to a^+}
\int_c^b f(x)\,dx
}
\]

when the limit exists.

Thus improper integration depends fundamentally on convergence theory.

%%%%%%%%%%%%%%%%%%%%%%%%%%%%%%%%%%%%%%%%%%%%%%%%%%%%%%%%%%%%
\subsection{Sequences of Functions}
\label{subsec:sequences-functions}
%%%%%%%%%%%%%%%%%%%%%%%%%%%%%%%%%%%%%%%%%%%%%%%%%%%%%%%%%%%%

A sequence of functions is written

\[
\boxed{
(f_n).
}
\]

Each

\[
f_n:A\to\R.
\]

There are several distinct notions of convergence for such sequences.

The two most fundamental are pointwise convergence and uniform convergence.

%%%%%%%%%%%%%%%%%%%%%%%%%%%%%%%%%%%%%%%%%%%%%%%%%%%%%%%%%%%%
\subsection{Pointwise Convergence}
\label{subsec:pointwise-convergence}
%%%%%%%%%%%%%%%%%%%%%%%%%%%%%%%%%%%%%%%%%%%%%%%%%%%%%%%%%%%%

\begin{definitionbox}{Pointwise Convergence}
A sequence \((f_n)\) converges pointwise to \(f\) on \(A\) if for every
fixed \(x\in A\),

\[
\boxed{
f_n(x)\to f(x).
}
\]
\end{definitionbox}

In logical form,

\[
\boxed{
\forall x\in A
\quad
\forall\varepsilon>0
\quad
\exists N=N(x,\varepsilon)
}
\]

such that

\[
n\geq N
\Longrightarrow
|f_n(x)-f(x)|<\varepsilon.
\]

The index \(N\) may depend on \(x\).

%%%%%%%%%%%%%%%%%%%%%%%%%%%%%%%%%%%%%%%%%%%%%%%%%%%%%%%%%%%%
\subsection{Uniform Convergence}
\label{subsec:uniform-convergence}
%%%%%%%%%%%%%%%%%%%%%%%%%%%%%%%%%%%%%%%%%%%%%%%%%%%%%%%%%%%%

\begin{definitionbox}{Uniform Convergence}
A sequence \((f_n)\) converges uniformly to \(f\) on \(A\) if for every
\(\varepsilon>0\), there exists \(N\) such that for every \(x\in A\),

\[
\boxed{
n\geq N
\Longrightarrow
|f_n(x)-f(x)|<\varepsilon.
}
\]
\end{definitionbox}

Equivalently,

\[
\boxed{
\sup_{x\in A}
|f_n(x)-f(x)|
\to0
}
\]

whenever the supremum formulation is appropriate.

The same \(N\) must work simultaneously for every point in the domain.

%%%%%%%%%%%%%%%%%%%%%%%%%%%%%%%%%%%%%%%%%%%%%%%%%%%%%%%%%%%%
\subsection{Pointwise Versus Uniform Convergence}
\label{subsec:pointwise-vs-uniform}
%%%%%%%%%%%%%%%%%%%%%%%%%%%%%%%%%%%%%%%%%%%%%%%%%%%%%%%%%%%%

Uniform convergence is stronger than pointwise convergence:

\[
\boxed{
\text{uniform convergence}
\Longrightarrow
\text{pointwise convergence}.
}
\]

The converse is false.

Consider

\[
f_n(x)=x^n
\]

on

\[
[0,1].
\]

Pointwise,

\[
f_n(x)\to
\begin{cases}
0, & 0\leq x<1,\\
1, & x=1.
\end{cases}
\]

The limiting function is discontinuous, even though every \(f_n\) is
continuous.

Therefore the convergence cannot be uniform.

%%%%%%%%%%%%%%%%%%%%%%%%%%%%%%%%%%%%%%%%%%%%%%%%%%%%%%%%%%%%
\subsection{Uniform Limits Preserve Continuity}
\label{subsec:uniform-limit-continuity}
%%%%%%%%%%%%%%%%%%%%%%%%%%%%%%%%%%%%%%%%%%%%%%%%%%%%%%%%%%%%

\begin{theorem}[Uniform Limit Theorem]
Suppose every \(f_n\) is continuous on \(A\) and

\[
f_n\to f
\]

uniformly on \(A\).

Then \(f\) is continuous on \(A\).
\end{theorem}

This is one of the central reasons uniform convergence is important.

Pointwise convergence alone does not preserve continuity.

%%%%%%%%%%%%%%%%%%%%%%%%%%%%%%%%%%%%%%%%%%%%%%%%%%%%%%%%%%%%
\subsection{Uniform Convergence and Integration}
\label{subsec:uniform-convergence-integration}
%%%%%%%%%%%%%%%%%%%%%%%%%%%%%%%%%%%%%%%%%%%%%%%%%%%%%%%%%%%%

\begin{theorem}
Suppose each \(f_n\) is Riemann integrable on \([a,b]\) and

\[
f_n\to f
\]

uniformly.

Then \(f\) is Riemann integrable and

\[
\boxed{
\lim_{n\to\infty}
\int_a^b f_n(x)\,dx
=
\int_a^b
\lim_{n\to\infty}f_n(x)\,dx.
}
\]
\end{theorem}

Thus uniform convergence allows the interchange

\[
\boxed{
\lim\int
=
\int\lim
}
\]

under these hypotheses.

%%%%%%%%%%%%%%%%%%%%%%%%%%%%%%%%%%%%%%%%%%%%%%%%%%%%%%%%%%%%
\subsection{Differentiation of Function Sequences}
\label{subsec:differentiation-function-sequences}
%%%%%%%%%%%%%%%%%%%%%%%%%%%%%%%%%%%%%%%%%%%%%%%%%%%%%%%%%%%%

Differentiation requires stronger hypotheses than integration.

A standard theorem states that if

\[
f_n(x_0)
\]

converges at some point \(x_0\), each \(f_n\) is continuously
differentiable, and

\[
f_n'\to g
\]

uniformly on a closed interval, then the sequence \((f_n)\) converges
uniformly to a differentiable function \(f\) satisfying

\[
\boxed{
f'=g.
}
\]

Thus one cannot generally interchange differentiation and limits using
pointwise convergence alone.

%%%%%%%%%%%%%%%%%%%%%%%%%%%%%%%%%%%%%%%%%%%%%%%%%%%%%%%%%%%%
\subsection{Infinite Series of Functions}
\label{subsec:function-series}
%%%%%%%%%%%%%%%%%%%%%%%%%%%%%%%%%%%%%%%%%%%%%%%%%%%%%%%%%%%%

A series of functions

\[
\sum_{n=1}^{\infty}f_n(x)
\]

is analyzed through its partial sums

\[
\boxed{
S_N(x)
=
\sum_{n=1}^{N}f_n(x).
}
\]

Pointwise or uniform convergence of the series means pointwise or uniform
convergence of the sequence

\[
(S_N).
\]

%%%%%%%%%%%%%%%%%%%%%%%%%%%%%%%%%%%%%%%%%%%%%%%%%%%%%%%%%%%%
\subsection{Weierstrass M-Test}
\label{subsec:weierstrass-m-test}
%%%%%%%%%%%%%%%%%%%%%%%%%%%%%%%%%%%%%%%%%%%%%%%%%%%%%%%%%%%%

\begin{theorem}[Weierstrass M-Test]
Suppose

\[
|f_n(x)|\leq M_n
\]

for every \(x\in A\), and suppose

\[
\sum_{n=1}^{\infty}M_n
\]

converges.

Then

\[
\boxed{
\sum_{n=1}^{\infty}f_n(x)
}
\]

converges uniformly and absolutely on \(A\).
\end{theorem}

The M-test converts a function-series problem into an ordinary numerical
series problem.

%%%%%%%%%%%%%%%%%%%%%%%%%%%%%%%%%%%%%%%%%%%%%%%%%%%%%%%%%%%%
\subsection{Power Series}
\label{subsec:power-series-analysis}
%%%%%%%%%%%%%%%%%%%%%%%%%%%%%%%%%%%%%%%%%%%%%%%%%%%%%%%%%%%%

A power series centered at \(a\) has the form

\[
\boxed{
\sum_{n=0}^{\infty}
c_n(x-a)^n.
}
\]

There exists a radius of convergence

\[
R\in[0,\infty]
\]

such that the series converges absolutely when

\[
|x-a|<R
\]

and diverges when

\[
|x-a|>R.
\]

The endpoints require separate analysis.

%%%%%%%%%%%%%%%%%%%%%%%%%%%%%%%%%%%%%%%%%%%%%%%%%%%%%%%%%%%%
\subsection{Uniform Convergence of Power Series}
\label{subsec:uniform-power-series}
%%%%%%%%%%%%%%%%%%%%%%%%%%%%%%%%%%%%%%%%%%%%%%%%%%%%%%%%%%%%

Inside the radius of convergence, power series behave especially well.

If

\[
0<r<R,
\]

then the power series converges uniformly on

\[
[a-r,a+r].
\]

Consequently, power series may be integrated and differentiated term by
term inside their interval of convergence.

%%%%%%%%%%%%%%%%%%%%%%%%%%%%%%%%%%%%%%%%%%%%%%%%%%%%%%%%%%%%
\subsection{Metric Spaces}
\label{subsec:metric-spaces-real-analysis}
%%%%%%%%%%%%%%%%%%%%%%%%%%%%%%%%%%%%%%%%%%%%%%%%%%%%%%%%%%%%

Many ideas from real analysis depend only on the notion of distance.

A metric space is a set \(X\) equipped with a distance function

\[
\boxed{
d:X\times X\to[0,\infty).
}
\]

The function \(d\) satisfies

\[
d(x,y)\geq0,
\]

\[
d(x,y)=0
\iff
x=y,
\]

\[
d(x,y)=d(y,x),
\]

and the triangle inequality

\[
\boxed{
d(x,z)
\leq
d(x,y)+d(y,z).
}
\]

%%%%%%%%%%%%%%%%%%%%%%%%%%%%%%%%%%%%%%%%%%%%%%%%%%%%%%%%%%%%
\subsection{Convergence in Metric Spaces}
\label{subsec:metric-space-convergence}
%%%%%%%%%%%%%%%%%%%%%%%%%%%%%%%%%%%%%%%%%%%%%%%%%%%%%%%%%%%%

A sequence

\[
(x_n)
\]

in a metric space \((X,d)\) converges to \(x\in X\) if

\[
\boxed{
d(x_n,x)\to0.
}
\]

Equivalently, for every \(\varepsilon>0\), there exists \(N\) such that

\[
n\geq N
\Longrightarrow
d(x_n,x)<\varepsilon.
\]

Thus ordinary real sequence convergence is the special case

\[
d(x,y)=|x-y|.
\]

%%%%%%%%%%%%%%%%%%%%%%%%%%%%%%%%%%%%%%%%%%%%%%%%%%%%%%%%%%%%
\subsection{Complete Metric Spaces}
\label{subsec:complete-metric-spaces}
%%%%%%%%%%%%%%%%%%%%%%%%%%%%%%%%%%%%%%%%%%%%%%%%%%%%%%%%%%%%

\begin{definitionbox}{Complete Metric Space}
A metric space is complete if every Cauchy sequence converges to a point
inside the space.
\end{definitionbox}

The real numbers are complete.

The rational numbers are not complete under the usual metric.

Completeness becomes fundamental in functional analysis and differential
equations.

%%%%%%%%%%%%%%%%%%%%%%%%%%%%%%%%%%%%%%%%%%%%%%%%%%%%%%%%%%%%
\subsection{Compactness in Metric Spaces}
\label{subsec:compactness-metric-analysis}
%%%%%%%%%%%%%%%%%%%%%%%%%%%%%%%%%%%%%%%%%%%%%%%%%%%%%%%%%%%%

For metric spaces, one useful characterization of compactness is:

\[
\boxed{
\text{every sequence has a convergent subsequence
whose limit lies in the space}.
}
\]

In \(\R^n\), the Heine--Borel Theorem states

\[
\boxed{
K\subseteq\R^n
\text{ is compact}
\iff
K\text{ is closed and bounded}.
}
\]

This equivalence does not hold in arbitrary metric spaces.

%%%%%%%%%%%%%%%%%%%%%%%%%%%%%%%%%%%%%%%%%%%%%%%%%%%%%%%%%%%%
\subsection{Connectedness and Intervals}
\label{subsec:connectedness-real-line}
%%%%%%%%%%%%%%%%%%%%%%%%%%%%%%%%%%%%%%%%%%%%%%%%%%%%%%%%%%%%

The connected subsets of \(\R\) are precisely the intervals.

This topological fact underlies the Intermediate Value Theorem.

A continuous image of a connected set is connected.

Thus if

\[
f:[a,b]\to\R
\]

is continuous, then

\[
f([a,b])
\]

must be an interval.

%%%%%%%%%%%%%%%%%%%%%%%%%%%%%%%%%%%%%%%%%%%%%%%%%%%%%%%%%%%%
\subsection{Contraction Mappings}
\label{subsec:contraction-mappings-analysis}
%%%%%%%%%%%%%%%%%%%%%%%%%%%%%%%%%%%%%%%%%%%%%%%%%%%%%%%%%%%%

\begin{definitionbox}{Contraction}
A map

\[
T:X\to X
\]

on a metric space is a contraction if there exists

\[
0\leq c<1
\]

such that

\[
\boxed{
d(Tx,Ty)
\leq
c\,d(x,y)
}
\]

for all \(x,y\in X\).
\end{definitionbox}

%%%%%%%%%%%%%%%%%%%%%%%%%%%%%%%%%%%%%%%%%%%%%%%%%%%%%%%%%%%%
\subsection{Banach Fixed-Point Theorem}
\label{subsec:banach-fixed-point}
%%%%%%%%%%%%%%%%%%%%%%%%%%%%%%%%%%%%%%%%%%%%%%%%%%%%%%%%%%%%

\begin{theorem}[Banach Fixed-Point Theorem]
Let \(X\) be a nonempty complete metric space and let

\[
T:X\to X
\]

be a contraction.

Then \(T\) has a unique fixed point

\[
x^*\in X
\]

satisfying

\[
\boxed{
T(x^*)=x^*.
}
\]

Moreover, repeated iteration

\[
x_{n+1}=T(x_n)
\]

converges to \(x^*\) for every initial point \(x_0\in X\).
\end{theorem}

This theorem is one of the most important bridges from real analysis to
differential equations and numerical mathematics.

%%%%%%%%%%%%%%%%%%%%%%%%%%%%%%%%%%%%%%%%%%%%%%%%%%%%%%%%%%%%
\subsection{Worked Example: Contraction}
\label{subsec:worked-contraction}
%%%%%%%%%%%%%%%%%%%%%%%%%%%%%%%%%%%%%%%%%%%%%%%%%%%%%%%%%%%%

\begin{workedexamplebox}{A Simple Fixed Point}

Consider

\[
T(x)=\frac{x+2}{3}
\]

on \(\R\).

For any \(x,y\),

\[
|T(x)-T(y)|
=
\left|
\frac{x-y}{3}
\right|
=
\frac13|x-y|.
\]

Thus \(T\) is a contraction with constant

\[
c=\frac13.
\]

A fixed point satisfies

\[
x=\frac{x+2}{3}.
\]

Therefore

\[
3x=x+2
\]

and

\[
x=1.
\]

\begin{answerbox}
\[
\boxed{
x^*=1.
}
\]
\end{answerbox}
\end{workedexamplebox}

%%%%%%%%%%%%%%%%%%%%%%%%%%%%%%%%%%%%%%%%%%%%%%%%%%%%%%%%%%%%
\subsection{Common Errors}
\label{subsec:real-analysis-common-errors}
%%%%%%%%%%%%%%%%%%%%%%%%%%%%%%%%%%%%%%%%%%%%%%%%%%%%%%%%%%%%

\subsubsection{Confusing Supremum with Maximum}

A supremum need not belong to the set.

A maximum must.

%%%%%%%%%%%%%%%%%%%%%%%%%%%%%%%%%%%%%%%%%%%%%%%%%%%%%%%%%%%%
\subsubsection{Treating Epsilon as a Fixed Small Number}

In a rigorous limit proof,

\[
\varepsilon>0
\]

is arbitrary.

The argument must work for every positive epsilon.

%%%%%%%%%%%%%%%%%%%%%%%%%%%%%%%%%%%%%%%%%%%%%%%%%%%%%%%%%%%%
\subsubsection{Choosing \(N\) Before Epsilon}

For sequence convergence, \(N\) is generally chosen in response to the
given \(\varepsilon\).

%%%%%%%%%%%%%%%%%%%%%%%%%%%%%%%%%%%%%%%%%%%%%%%%%%%%%%%%%%%%
\subsubsection{Reversing Quantifiers}

The statement

\[
\forall\varepsilon>0\;\exists N
\]

is fundamentally different from

\[
\exists N\;\forall\varepsilon>0.
\]

Quantifier order matters.

%%%%%%%%%%%%%%%%%%%%%%%%%%%%%%%%%%%%%%%%%%%%%%%%%%%%%%%%%%%%
\subsubsection{Assuming Bounded Means Convergent}

Bounded sequences need not converge.

For example,

\[
(-1)^n
\]

is bounded and divergent.

%%%%%%%%%%%%%%%%%%%%%%%%%%%%%%%%%%%%%%%%%%%%%%%%%%%%%%%%%%%%
\subsubsection{Assuming \(a_n\to0\) Makes a Series Converge}

The condition

\[
a_n\to0
\]

is necessary but not sufficient.

%%%%%%%%%%%%%%%%%%%%%%%%%%%%%%%%%%%%%%%%%%%%%%%%%%%%%%%%%%%%
\subsubsection{Confusing Cauchy with Convergent in Arbitrary Spaces}

Every convergent sequence is Cauchy.

A Cauchy sequence converges only when the ambient metric space is complete.

%%%%%%%%%%%%%%%%%%%%%%%%%%%%%%%%%%%%%%%%%%%%%%%%%%%%%%%%%%%%
\subsubsection{Assuming Continuity Implies Uniform Continuity}

Continuity on an arbitrary domain does not imply uniform continuity.

Compactness often supplies the missing condition.

%%%%%%%%%%%%%%%%%%%%%%%%%%%%%%%%%%%%%%%%%%%%%%%%%%%%%%%%%%%%
\subsubsection{Assuming Continuity Implies Differentiability}

Differentiability implies continuity, but continuity does not imply
differentiability.

%%%%%%%%%%%%%%%%%%%%%%%%%%%%%%%%%%%%%%%%%%%%%%%%%%%%%%%%%%%%
\subsubsection{Assuming Every Bounded Function Is Riemann Integrable}

Boundedness is necessary in the standard Riemann theory on closed
intervals, but it is not sufficient.

%%%%%%%%%%%%%%%%%%%%%%%%%%%%%%%%%%%%%%%%%%%%%%%%%%%%%%%%%%%%
\subsubsection{Confusing Pointwise and Uniform Convergence}

Uniform convergence requires a single index \(N\) that works for every
point in the domain.

Pointwise convergence allows \(N\) to depend on the point.

%%%%%%%%%%%%%%%%%%%%%%%%%%%%%%%%%%%%%%%%%%%%%%%%%%%%%%%%%%%%
\subsubsection{Interchanging Limits Without Justification}

Expressions such as

\[
\lim\int f_n
\]

and

\[
\int\lim f_n
\]

are not automatically equal.

Similarly,

\[
\lim f_n'
\]

need not equal

\[
\left(\lim f_n\right)'.
\]

Specific hypotheses are required.

%%%%%%%%%%%%%%%%%%%%%%%%%%%%%%%%%%%%%%%%%%%%%%%%%%%%%%%%%%%%
\subsection{Applications}
\label{subsec:real-analysis-applications}
%%%%%%%%%%%%%%%%%%%%%%%%%%%%%%%%%%%%%%%%%%%%%%%%%%%%%%%%%%%%

\begin{applicationbox}
\textbf{Calculus Foundations.}
Real analysis rigorously explains limits, continuity, derivatives,
integrals, and infinite series.

\medskip

\textbf{Differential Equations.}
Existence and uniqueness results frequently rely on completeness,
contraction mappings, compactness, and convergence arguments.

\medskip

\textbf{Numerical Analysis.}
Convergence of algorithms requires rigorous control of approximation error
and limiting behavior.

\medskip

\textbf{Probability.}
Random variables, expectations, probability measures, and convergence
theorems build upon analytical ideas.

\medskip

\textbf{Optimization.}
Continuity and compactness provide existence of extrema, while convexity
and differentiability provide structural information.

\medskip

\textbf{Functional Analysis.}
Sequences and completeness are extended from numbers to infinite-dimensional
spaces of functions and operators.

\medskip

\textbf{Topology.}
Continuity, compactness, connectedness, and metric spaces connect analysis
directly to point-set topology.

\medskip

\textbf{Measure Theory.}
The limitations of Riemann integration motivate the development of
Lebesgue measure and integration.
\end{applicationbox}

%%%%%%%%%%%%%%%%%%%%%%%%%%%%%%%%%%%%%%%%%%%%%%%%%%%%%%%%%%%%
\subsection{Exercises}
\label{subsec:real-analysis-exercises}
%%%%%%%%%%%%%%%%%%%%%%%%%%%%%%%%%%%%%%%%%%%%%%%%%%%%%%%%%%%%

\begin{exercise}[1]
Find the supremum and infimum of

\[
A=(2,7).
\]

Does \(A\) have a maximum or minimum?
\end{exercise}

\begin{exercise}[1]
Determine whether

\[
x_n=\frac{n}{n+1}
\]

is bounded and monotone.
\end{exercise}

\begin{exercise}[1]
Determine the limit of

\[
x_n=\frac{3n+1}{2n+5}.
\]
\end{exercise}

\begin{exercise}[1]
Determine whether

\[
(-1)^n
\]

converges.
\end{exercise}

\begin{exercise}[1]
Determine whether

\[
\sum_{n=1}^{\infty}\frac1{2^n}
\]

converges.
\end{exercise}

\begin{exercise}[2]
Using the epsilon definition, prove

\[
\frac{2}{n}\to0.
\]
\end{exercise}

\begin{exercise}[2]
Prove that every convergent sequence is bounded.
\end{exercise}

\begin{exercise}[2]
Prove that the limit of a convergent sequence is unique.
\end{exercise}

\begin{exercise}[2]
Find

\[
\limsup_{n\to\infty}(-1)^n
\]

and

\[
\liminf_{n\to\infty}(-1)^n.
\]
\end{exercise}

\begin{exercise}[2]
Show that every Cauchy sequence is bounded.
\end{exercise}

\begin{exercise}[2]
Use an epsilon--delta argument to prove

\[
\lim_{x\to3}2x=6.
\]
\end{exercise}

\begin{exercise}[2]
Prove directly from the definition that

\[
f(x)=5x-2
\]

is continuous at every real number.
\end{exercise}

\begin{exercise}[2]
Show that

\[
f(x)=x^2
\]

is not uniformly continuous on \(\R\).
\end{exercise}

\begin{exercise}[2]
Show that

\[
f(x)=x^2
\]

is uniformly continuous on

\[
[0,1].
\]
\end{exercise}

\begin{exercise}[2]
Use the Intermediate Value Theorem to prove that

\[
x^3+x-1=0
\]

has a solution in \((0,1)\).
\end{exercise}

\begin{exercise}[2]
Determine whether the function

\[
f(x)=
\begin{cases}
0, & x<0,\\
1, & x\geq0
\end{cases}
\]

is Riemann integrable on \([-1,1]\).
\end{exercise}

\begin{exercise}[3]
Prove the Monotone Convergence Theorem for increasing bounded sequences
using the least upper bound property.
\end{exercise}

\begin{exercise}[3]
Prove that

\[
x_n=\left(1+\frac1n\right)^n
\]

is bounded and investigate its monotonicity.
\end{exercise}

\begin{exercise}[3]
Show that a sequence converges to \(L\) if and only if every subsequence
has a further subsequence converging to \(L\).
\end{exercise}

\begin{exercise}[3]
Prove the sequential characterization of continuity.
\end{exercise}

\begin{exercise}[3]
Use the Mean Value Theorem to prove that if

\[
f'(x)=0
\]

on an interval, then \(f\) is constant there.
\end{exercise}

\begin{exercise}[3]
Suppose

\[
|f'(x)|\leq M
\]

for every \(x\in(a,b)\). Prove

\[
|f(x)-f(y)|
\leq
M|x-y|.
\]
\end{exercise}

\begin{exercise}[3]
Show that every Lipschitz function is uniformly continuous.
\end{exercise}

\begin{exercise}[3]
Determine the pointwise limit of

\[
f_n(x)=x^n
\]

on \([0,1]\), and prove that the convergence is not uniform.
\end{exercise}

\begin{exercise}[3]
Show that

\[
f_n(x)=\frac{x}{n}
\]

converges uniformly to \(0\) on \([0,1]\).
\end{exercise}

\begin{exercise}[3]
Use the Weierstrass M-test to show that

\[
\sum_{n=1}^{\infty}
\frac{x^n}{n^2}
\]

converges uniformly on \([-1,1]\).
\end{exercise}

\begin{exercise}[3]
Determine the radius of convergence of

\[
\sum_{n=1}^{\infty}
\frac{x^n}{n}.
\]
\end{exercise}

\begin{exercise}[4]
Prove the Bolzano--Weierstrass Theorem for bounded sequences in \(\R\).
\end{exercise}

\begin{exercise}[4]
Prove that every Cauchy sequence of real numbers converges using
completeness.
\end{exercise}

\begin{exercise}[4]
Prove the Heine--Cantor Theorem for continuous functions on compact metric
spaces.
\end{exercise}

\begin{exercise}[4]
Prove that a continuous image of a compact set is compact.
\end{exercise}

\begin{exercise}[4]
Prove that a continuous image of a connected set is connected.
\end{exercise}

\begin{exercise}[4]
Prove that the uniform limit of continuous functions is continuous.
\end{exercise}

\begin{exercise}[4]
Prove that uniform convergence permits interchange of limit and Riemann
integration on a closed bounded interval.
\end{exercise}

\begin{exercise}[4]
Construct a sequence of differentiable functions that converges uniformly
while the sequence of derivatives does not converge uniformly.
\end{exercise}

\begin{exercise}[4]
Prove the Banach Fixed-Point Theorem.
\end{exercise}

%%%%%%%%%%%%%%%%%%%%%%%%%%%%%%%%%%%%%%%%%%%%%%%%%%%%%%%%%%%%
\subsection{Conceptual Summary}
\label{subsec:real-analysis-summary}
%%%%%%%%%%%%%%%%%%%%%%%%%%%%%%%%%%%%%%%%%%%%%%%%%%%%%%%%%%%%

Real analysis provides the rigorous foundations of calculus.

The completeness of \(\R\) guarantees the existence of least upper bounds:

\[
\boxed{
A\neq\varnothing,\quad
A\text{ bounded above}
\Longrightarrow
\sup A\in\R.
}
\]

Sequence convergence is defined by

\[
\boxed{
\forall\varepsilon>0
\quad
\exists N
\quad
n\geq N
\Rightarrow
|x_n-L|<\varepsilon.
}
\]

Completeness also gives the Cauchy criterion:

\[
\boxed{
(x_n)\text{ converges in }\R
\iff
(x_n)\text{ is Cauchy}.
}
\]

Function limits are defined by

\[
\boxed{
\forall\varepsilon>0
\quad
\exists\delta>0
\quad
0<|x-a|<\delta
\Rightarrow
|f(x)-L|<\varepsilon.
}
\]

Continuity means

\[
\boxed{
\lim_{x\to a}f(x)=f(a).
}
\]

Compactness strengthens continuity through results such as the Extreme
Value and Heine--Cantor Theorems.

Differentiability is a stronger local property than continuity:

\[
\boxed{
\text{differentiable}
\Longrightarrow
\text{continuous}.
}
\]

The Mean Value Theorem connects derivatives to global changes:

\[
\boxed{
f'(c)
=
\frac{f(b)-f(a)}{b-a}.
}
\]

Riemann integration is constructed using upper and lower approximations:

\[
\boxed{
U(f,P)-L(f,P)\to0.
}
\]

For sequences of functions,

\[
\boxed{
\text{uniform convergence}
\Longrightarrow
\text{pointwise convergence},
}
\]

but not conversely.

Uniform convergence preserves important analytical structure, including
continuity and, under appropriate hypotheses, integration.

Metric spaces reveal that many analytical ideas depend only on distance.

Completeness becomes

\[
\boxed{
\text{every Cauchy sequence converges}.
}
\]

Together, these ideas establish the rigorous language needed for the deeper
study of measure, integration, probability, functional analysis, harmonic
analysis, differential equations, and modern analysis.

%%%%%%%%%%%%%%%%%%%%%%%%%%%%%%%%%%%%%%%%%%%%%%%%%%%%%%%%%%%%
\subsection{Section Connections}
\label{subsec:real-analysis-connections}
%%%%%%%%%%%%%%%%%%%%%%%%%%%%%%%%%%%%%%%%%%%%%%%%%%%%%%%%%%%%

\chapterconnection{
Real Analysis places the computational ideas of Single-Variable and
Multivariable Calculus on a rigorous foundation. Completeness explains why
limits exist in the real number system, while sequences and Cauchy
sequences provide the basic language of convergence. Continuity,
compactness, connectedness, differentiation, and Riemann integration link
analysis directly to Point-Set Topology. Uniform convergence becomes
essential when limits are taken through operations such as integration and
differentiation. The limitations of the Riemann integral motivate Measure
Theory and Lebesgue Integration, while complete metric spaces and
contraction mappings lead toward Functional Analysis. These ideas also
provide the theoretical foundations for Probability, Differential
Equations, Numerical Analysis, Optimization, Fourier Analysis, and
Harmonic Analysis.
}

%%%%%%%%%%%%%%%%%%%%%%%%%%%%%%%%%%%%%%%%%%%%%%%%%%%%%%%%%%%%
% SECTION SOURCE TRACKING
%%%%%%%%%%%%%%%%%%%%%%%%%%%%%%%%%%%%%%%%%%%%%%%%%%%%%%%%%%%%

%------------------------------------------------------------
% 5.5 Real Analysis
%------------------------------------------------------------
%
% Source basis:
%   - Standard undergraduate real analysis.
%   - Standard rigorous development of the real number system,
%     sequences, continuity, differentiation, Riemann integration,
%     and sequences of functions.
%
% Used for:
%   - completeness of the real numbers;
%   - upper and lower bounds;
%   - supremum and infimum;
%   - Archimedean property;
%   - density of rational numbers;
%   - sequences and convergence;
%   - uniqueness and algebra of limits;
%   - bounded and monotone sequences;
%   - Monotone Convergence Theorem for sequences;
%   - subsequences;
%   - Bolzano--Weierstrass Theorem;
%   - limsup and liminf;
%   - Cauchy sequences;
%   - series and absolute convergence;
%   - epsilon--delta limits;
%   - sequential criteria;
%   - continuity;
%   - Intermediate Value Theorem;
%   - Extreme Value Theorem;
%   - uniform continuity;
%   - Heine--Cantor Theorem;
%   - differentiability;
%   - Rolle's Theorem;
%   - Mean Value Theorem;
%   - Lipschitz continuity;
%   - Taylor's Theorem;
%   - Riemann integration;
%   - upper and lower sums;
%   - improper integrals;
%   - pointwise convergence;
%   - uniform convergence;
%   - sequences and series of functions;
%   - Weierstrass M-test;
%   - power series;
%   - metric spaces;
%   - completeness and compactness;
%   - contraction mappings;
%   - Banach Fixed-Point Theorem.
%
% Notes:
%   - Elementary computational limit, derivative, integral, and
%     series techniques are developed earlier in Calculus.
%   - General point-set topology is developed canonically in
%     Chapter 4.
%   - Measure-theoretic convergence and integration are developed
%     in the following analysis sections.
%   - Infinite-dimensional complete spaces are developed more fully
%     in Functional Analysis.
%
%%%%%%%%%%%%%%%%%%%%%%%%%%%%%%%%%%%%%%%%%%%%%%%%%%%%%%%%%%%%